\documentclass[11pt, a4paper]{article}
\usepackage[utf8]{inputenc}
\usepackage[spanish]{babel}
\usepackage{geometry}
\geometry{margin=2.5cm}
\usepackage{tabularx}
\usepackage{booktabs}
\usepackage{hyperref}

\title{\textbf{Código de Ética y Reglamento Interno de Grupo} \cite{3, 4}}
\author{MAT1186 Introducción al Cálculo \cite{13}}
\date{Universidad Católica de Temuco \cite{15}}

\begin{document}

\maketitle

\section*{Preámbulo}
El presente documento regula el funcionamiento del grupo de trabajo para el desarrollo de la Evaluación Integrada de Desempeño N°1 \cite{16}. Este surge del compromiso colectivo con los valores de responsabilidad, honestidad académica, colaboración activa y respeto mutuo \cite{17}. Asumimos que el éxito del proyecto depende de la participación equitativa \cite{18}. Al suscribir este código, cada integrante acepta voluntariamente su cumplimiento \cite{19}.

\section{Identificación del Grupo}
\begin{table}[h!]
\centering
\begin{tabularx}{\textwidth}{lX}
\toprule
\textbf{Rol} & \textbf{Nombre} \\
\midrule
Líder & Benjamín De La Fuente \cite{24} \\
Suplente & Sebastián Rodrigo Rivera Nahuelhuen \cite{24} \\
Integrante & Rodrigo Reyes \cite{24} \\
\bottomrule
\end{tabularx}
\end{table}

\section{Definición de Roles y Responsabilidades}

\subsection{Líder de Grupo}
\begin{itemize}
    \item Organizar y convocar las reuniones de trabajo, fijando agenda y plazos internos \cite{27}.
    \item Ser el canal oficial de comunicación con el equipo docente \cite{28}.
    \item Supervisar el avance general del proyecto y velar por los plazos \cite{29}.
    \item Gestionar el repositorio de GitHub y revisar los commits de los integrantes \cite{30}.
    \item Mediar en conflictos internos antes de escalarlos al docente \cite{31}.
    \item Participar activamente en el desarrollo técnico del proyecto \cite{32}.
\end{itemize}

\subsection{Suplente}
\begin{itemize}
    \item Asumir las funciones de líder en ausencia justificada de este \cite{34}.
    \item Desarrollar los módulos técnicos asignados en coordinación con el equipo \cite{35}.
    \item Registrar avances mediante commits regulares y descriptivos en GitHub \cite{36}.
    \item Participar activamente en la preparación de la defensa oral \cite{37}.
    \item Informar oportunamente al líder si alguna tarea no puede completarse \cite{38}.
\end{itemize}

\subsection{Integrante}
\begin{itemize}
    \item Desarrollar los módulos técnicos asignados según la distribución acordada \cite{43}.
    \item Registrar avances mediante commits regulares y descriptivos en GitHub \cite{44}.
    \item Participar activamente en la preparación de la defensa oral \cite{45}.
    \item Documentar el código bajo los estándares acordados por el grupo \cite{46}.
    \item Informar oportunamente al líder si alguna tarea no puede completarse \cite{47}.
\end{itemize}

\section{Funcionamiento Interno}

\subsection{Reuniones y comunicación}
\begin{itemize}
    \item El grupo se reunirá con una frecuencia mínima de dos veces por semana \cite{51}.
    \item El canal principal de comunicación será WhatsApp, respondiendo en máximo 24 horas \cite{53}.
    \item La inasistencia a una reunión deberá ser comunicada con 12 horas de anticipación \cite{55}.
\end{itemize}

\subsection{Distribución de tareas}
\begin{itemize}
    \item Las tareas se distribuirán en la primera reunión formal buscando equidad \cite{58}.
    \item Cada integrante es responsable de completar sus tareas en los plazos acordados \cite{60}.
\end{itemize}

\subsection{Repositorio y control de versiones}
\begin{itemize}
    \item Todos deben realizar un mínimo de tres commits significativos en GitHub \cite{67}.
    \item Queda prohibido realizar commits con trabajo ajeno como si fuera propio \cite{70}.
\end{itemize}

\section{Código de Ética y Conducta}

\subsection{Compromisos individuales}
\begin{itemize}
    \item \textbf{Uso de IA:} Se permite apoyo de IA, pero todos deben comprender y defender el trabajo resultante \cite{75, 76}.
    \item \textbf{Participación:} Ningún integrante será solo observador; todos deben contribuir \cite{79, 80}.
    \item \textbf{Transparencia:} Si no se comprende un módulo, debe comunicarse al grupo para el aprendizaje mutuo \cite{81, 82}.
    \item \textbf{Respeto:} Interacciones con tolerancia, prohibiendo descalificaciones y discriminación \cite{83, 84}.
    \item \textbf{Puntualidad:} Cumplir plazos internos y asistir puntualmente a instancias de trabajo \cite{85}.
    \item \textbf{Responsabilidad:} Todos deben comprender el funcionamiento completo del programa \cite{86}.
    \item \textbf{Integridad:} El historial de commits debe reflejar el trabajo real; no manipularlo \cite{91, 92}.
\end{itemize}

\subsection{Consecuencias por incumplimiento}
\begin{table}[h!]
\centering
\begin{tabularx}{\textwidth}{clX}
\toprule
\textbf{Nivel} & \textbf{Tipo de falta} & \textbf{Medida} \\
\midrule
1 & Leve (primera vez) & Llamado de atención verbal por parte del líder \cite{95}. \\
2 & Moderada (reincidencia) & Reasignación de tareas y advertencia formal \cite{95}. \\
3 & Grave & Comunicación formal al equipo docente solicitando intervención \cite{95}. \\
\bottomrule
\end{tabularx}
\end{table}

\section{Procedimiento ante Conflictos}
\begin{itemize}
    \item \textbf{Diálogo directo:} Intentar resolver el conflicto autónomamente en 48 horas \cite{98}.
    \item \textbf{Mediación:} Si no hay acuerdo, el líder o suplente actuará como mediador \cite{99, 100}.
    \item \textbf{Escalamiento:} Si persiste, comunicar la situación al docente con registro escrito \cite{101}.
    \item Las decisiones estructurales deben ser consensuadas por mayoría (2 de 3) \cite{102}.
\end{itemize}

\section{Declaración Final}
Los integrantes asumen el compromiso de actuar conforme a los principios detallados, reconociendo que su incumplimiento afectará la calificación grupal y la integridad de su aprendizaje profesional \cite{108, 109}.

\end{document}